\begin{helper}[Approximation of the degenerate case]
Let $n\ge2$, $p>1$ and $r>0$.  Then there is $r'=r'(n,p,r)>0$ with the following
property.  Let $\gamma_0>0$ and $\gamma_1\le r\gamma_0$, let
$A\colon\R^n\to\R^n$ satisfy $A(0)=0$ and (Q1) of the thesis--(Q2) of the thesis with
$\varepsilon=0$ and constants $\gamma_0,\gamma_1$
(\noderef{EllipticStructureQ12}), and assume the cited solvability result (C2)
(\noderef{CitedDirichletW1p}) \emph{for every} field
$A'\colon\R^n\to\R^n$ and all parameters $\varepsilon,\gamma_0',\gamma_1'$ with
$\gamma_0'>0$ and $\gamma_1'>0$.  Let
$B_R=B_R(x_0)$, $R>0$, let $u\in W^{1,p}(B_R)$ with weak gradient $g=\nabla u$
(\noderef{MemW1pOn}), and let $u$
be a weak solution of $\operatorname{div}A(\nabla u)=0$ in $B_R$
(\noderef{IsWeakSolutionOn}) with the same weak gradient $g$.  Then there exist
a sequence $e_m>0$ with $e_m\to0$, vector fields $A_{e_m}$ and functions
$u_m$ with weak gradients $g_m$ such that
\begin{enumerate}
\item for every $m$, $A_{e_m}(0)=0$ and there are constants
  $\gamma_0^{(m)}>0$ and $\gamma_1^{(m)}\le r'\gamma_0^{(m)}$ for which
  $A_{e_m}$ satisfies (Q1) of the thesis--(Q2) of the thesis with parameter
  $\varepsilon=e_m$ and those constants (\noderef{EllipticStructureQ12});
\item $u_m$ is a weak solution of $\operatorname{div}A_{e_m}(\nabla u_m)=0$ in
  $B_R$, for every $m$ (\noderef{IsWeakSolutionOn});
\item $u_m-u\in W^{1,p}_0(B_R)$ for every $m$
  (\noderef{MemW1pZeroOn}), i.e.\ $u_m$ has the same boundary values as $u$;
\item $\int_{B_R}|g_m-g|^p\,dx\to0$ as $m\to\infty$; and
\item $g_m(x)\to g(x)$ as $m\to\infty$, for almost every $x\in B_R$.
\end{enumerate}

\textbf{On the two quantifiers.}  Clause (1) asks only that the \emph{ratio} of
the approximating structure constants be bounded, by a number $r'$ chosen before
the family, and not that the constants themselves be $\gamma_0,\gamma_1$.  The
stronger form is false for the natural perturbation: for $p=6$ and
$A(h)=|h|^4h$, which satisfies (Q1) of the thesis--(Q2) of the thesis at
$\varepsilon=0$ with $\gamma_0=1$ and $\gamma_1=5$ because
$(D A(h)\xi,\xi)=|h|^4|\xi|^2+4|h|^2(h,\xi)^2$, the field
$A_e(h)=A(h)+e^{\frac p2-1}h=A(h)+e^2h$ has, at $|h|^2=e$ and $\xi\perp h$ with
$|\xi|=1$,
\[
(DA_e(h)\xi,\xi)=e^2+e^2=2e^2,
\qquad
\gamma_0\bigl(e+|h|^2\bigr)^{\frac p2-1}|\xi|^2=(2e)^2=4e^2 ,
\]
so (Q1) of the thesis fails at $\gamma_0=1$; the largest admissible constant there is
$\tfrac12$.  The bounded-ratio form is what the consumers need, because the
constants of Theorems~Theorem~1 of the thesis and~Theorem~2 of the thesis depend only on $n$, $p$ and
$\gamma_1/\gamma_0$, and both consumers bind their structure constants through
an inequality $\gamma_1\le r\gamma_0$: they instantiate that bound at $r'$.

Clauses (4) and (5) are stated along a \emph{sequence} $e_m\to0$ and not for the
whole family $(A_e)_{e>0}$.  That is the strongest form available: convergence
in $L^p$ yields almost everywhere convergence only along a sequence, by passage
to a subsequence, and never for a family indexed by a continuum.

\textbf{On the form of the (C2) hypothesis.}  The cited result is assumed for
every field and every parameter, not only for the given $A$ at
$\varepsilon=0$, because the solutions $u_m$ whose existence is asserted solve
the equation of the \emph{regularized} field $A_{e_m}$, and those fields are
produced inside the conclusion, so no hypothesis attached to a single named
field could reach them.  This is the form (C2) has in the paper, which states it
for an arbitrary operator: ``If $A$ satisfies (Q1) of the thesis--(Q2) of the thesis with
some $\varepsilon\ge0$ and $A(0)=0$, $B$ is a ball and $g\in W^{1,p}(B)$, then
there is a weak solution $v$ of (3.1) of the thesis in $B$ with
$v-g\in W^{1,p}_0(B)$'' (Note for the formalization, (C2)).  So the universally
quantified clause is the faithful transcription and the earlier
single-instance one was the weakening.  The quantifier is restricted to
$\gamma_0'>0$ and $\gamma_1'>0$, and those side conditions are not optional: the
structure conditions carry positive constants in the thesis
(``$\gamma_0,\gamma_1>0$''), so a clause quantified over $\gamma_0'\le0$ or
$\gamma_1'\le0$ would assert more than the cited result does --- and since this
is a \emph{hypothesis} of the present node, assuming more than the literature
gives risks making the node vacuously true.  The remaining antecedents need no
guard, because the predicate \noderef{CitedDirichletW1p} already carries
(Q1) of the thesis--(Q2) of the thesis, $\varepsilon\ge0$ and $A'(0)=0$ inside its own
statement.

\textbf{On the hypothesis $u\in W^{1,p}(B_R)$.}  The membership of $u$ in
$W^{1,p}$ of the \emph{whole} ball, and not merely in
$W^{1,p}_{\mathrm{loc}}(B_R)$, is a restoration of what the cited result the
conclusion rests on actually requires, not a strengthening of this node.  The
cited solvability result (C2) reads: ``If $A$ satisfies
(Q1) of the thesis--(Q2) of the thesis with some $\varepsilon\ge0$ and $A(0)=0$, $B$ is a
ball and $g\in W^{1,p}(B)$, then there is a weak solution $v$ of
(3.1) of the thesis in $B$ with $v-g\in W^{1,p}_0(B)$'' (Note for the formalization,
(C2)); the boundary datum is required to lie in $W^{1,p}$ of the ball.  The
conclusion of the present node applies (C2) with the boundary datum $u$ --- that
is how the approximating solutions $u_m$ are produced and how
$u_m-u\in W^{1,p}_0(B_R)$ is asserted in clause (3) --- so without this
hypothesis clauses (2)--(5) would assert more than (C2) licenses.  The local
membership that \noderef{IsWeakSolutionOn} carries does not supply it: a
$W^{1,p}_{\mathrm{loc}}(B_R)$ function need not be in $L^p(B_R)$, because $B_R$
is not compactly contained in itself.  Nor does the hypothesis narrow the node's
use: both consumers apply it on a ball whose \emph{closed} ball lies inside the
solution's domain, and for such a ball the local membership on the domain does
give the global one on the ball.

\textbf{Why a bounded ratio is enough.}  The paper itself records that ``it is
enough to prove Theorems~Theorem~1 of the thesis and~Theorem~2 of the thesis when $\varepsilon>0$,
since our constants are independent of $\varepsilon$'', and recovers
$\varepsilon=0$ ``by an approximation argument as done in [DB, section~2]''
(thesis p.~16, the remark after Theorem~Theorem~2 of the thesis).  Since those constants
depend on the structure only through $\gamma_1/\gamma_0$, a bound on that ratio
holding uniformly along the sequence is exactly what the reduction consumes, and
the identical-constants clause was never needed.  Both consumers
use exactly the sequence form.
\end{helper}
\begin{proof}
This is the approximation argument of [DB, section 2], which the paper's Note
for the formalization requires to be proved.  The proof constructs the
approximating family and its structure constants (clause (1)), produces the
regularized solutions from the cited result (C2) (clauses (2) and (3)), proves
the convergence of their gradients in $L^p(B_R)$ (clause (4)), and extracts the
subsequence along which the gradients also converge almost everywhere (clause
(5)).

\emph{Reading the structure conditions pointwise.}  This is recorded first
because every later paragraph uses it, including the construction of the family.
Both (Q1) of the thesis and (Q2) of the thesis of \noderef{EllipticStructureQ12} are
almost-everywhere statements in $x$, and the fields appearing here do not depend
on $x$: the hypothesis of the theorem is
$\mathrm{EllipticStructureQ12}(n,p,0,\gamma_0,\gamma_1,\R^n,(x,z,h)\mapsto A(h))$,
and the same shape will be produced for $A_e$.  For such a field the property
quantified over $x$ does not involve $x$, so the exceptional set is either empty
or all of $\R^n$; the latter is impossible because $\R^n$ is not Lebesgue null,
containing a ball of positive volume $\omega_n\rho^n>0$ (Preamble: the volume of
a ball).  Hence for an $x$-independent field (Q1) of the thesis and (Q2) of the thesis
hold \emph{for every} $x$, i.e.\ as statements about the field of $h$ alone, and
we use them in that form throughout, for $A$ from here on and for $A_e$ as soon
as it has been shown to satisfy the bundle.  The same remark applies to the two
cited nodes used below, \noderef{StructureMonotonicityBound} and
\noderef{StructureGrowthBound}, whose conclusions are likewise ``for a.e.\ $x$''
with $x$-free right-hand sides.  There is no circularity: the argument uses only
that the field does not depend on $x$ and that $\R^n$ is not null.

\emph{The approximating family.}  Fix once and for all the mollifier
$\rho=\rho_1$ of the Preamble, with the two further properties recorded there:
$\rho\in C^\infty(\R^n)$, $\rho\ge0$, $\operatorname{supp}\rho\subseteq\overline{B_1(0)}$,
$\int\rho=1$, $\rho\equiv\kappa>0$ on $\overline{B_{7/8}(0)}$ and
$\|\rho\|_\infty<\infty$.  For $e>0$ put
\[
A_e(h)=\int_{\R^n}A(h+\sqrt e\,z)\,\rho(z)\,dz
-\int_{\R^n}A(\sqrt e\,z)\,\rho(z)\,dz ,
\]
that is, $A$ mollified at scale $\sqrt e$ and renormalised so that $A_e(0)=0$.
The subtracted vector is a constant, so it cancels in every difference
$A_e(a)-A_e(b)$ and does not affect the derivative.

\emph{$A_e$ is a convolution, hence $C^\infty$.}  Put
\[
\sigma(w)=e^{-n/2}\rho\Bigl(-\frac w{\sqrt e}\Bigr).
\]
Then $\sigma\in C^\infty(\R^n)$ with
$\operatorname{supp}\sigma\subseteq\overline{B_{\sqrt e}(0)}$ compact, both
because $\rho$ has those properties and $w\mapsto-w/\sqrt e$ is a smooth
diffeomorphism.  For every $h$, the change of variables $w=-\sqrt e\,z$ --- a
unimodular linear map composed with a dilation, so $dw=e^{n/2}dz$ (Preamble:
invariance of Lebesgue measure under translations and unimodular linear maps,
together with the scaling in the volume-of-a-ball item) --- gives
\[
(A*\sigma)(h)=\int_{\R^n}A(h-w)\,\sigma(w)\,dw
=\int_{\R^n}A(h+\sqrt e\,z)\,e^{-n/2}\rho(z)\,e^{n/2}\,dz
=\int_{\R^n}A(h+\sqrt e\,z)\,\rho(z)\,dz .
\]
Hence $A_e=A*\sigma-(A*\sigma)(0)$, componentwise.  Now $A$ is continuous on all
of $\R^n$ by clause \texttt{hcont} of \noderef{EllipticStructureQ12}, so each
$A^i$ is locally bounded and therefore in $L^1_{\mathrm{loc}}(\R^n)$ (a
continuous function on a compact set is bounded there, and a bounded measurable
function on a set of finite measure is integrable --- Preamble, clause (a) of the
domination item, against the constant majorant).  The Preamble item on
convolution with a smooth compactly supported kernel therefore applies to each
component with $w=A^i$ and the kernel $\sigma$, and gives
\[
A_e\in C^\infty(\R^n),
\qquad
D_j A_e^i=A^i*(D_j\sigma)\qquad(i,j=1,\dots,n).
\]
\emph{This is the only route to the smoothness of $A_e$ used here.}  The formula
$DA_e(h)=\int DA(h+\sqrt e\,z)\rho(z)\,dz$ --- differentiation under the integral
sign --- is \emph{not} used and is not available when $1<p<2$: the only
pointwise majorant for $|DA|$ that (Q2) of the thesis supplies at $\varepsilon=0$ is a
constant multiple of $|h'|^{p-2}$, and a majorant dominating the $h$-derivative
uniformly for $h$ near $0$ would have to be $+\infty$ on a set of positive
measure.  What replaces it is the increment route of the next two paragraphs,
which differentiates nothing under an integral sign.  Finally
$A_e\to A$ locally uniformly as $e\downarrow0$, because $A$ is continuous and
$\int\rho=1$.

Both structure conditions for $A_e$ now rest on one two-sided estimate for the
scalar kernel --- its upper half for (Q2) of the thesis, its lower half for
(Q1) of the thesis when $p\ge2$, the case $1<p<2$ of (Q1) of the thesis needing only the
trivial bound $|h+\sqrt e\,z|^2\le2(e+|h|^2)$ ---
\[
K(h,e)=\int_{\R^n}\bigl(|h+\sqrt e\,z|^2\bigr)^{\frac p2-1}\rho(z)\,dz ,
\]
the \emph{weight sandwich}: there are constants $0<c_1\le C_2<\infty$ depending
only on $n$ and $p$ with
\[
c_1\bigl(e+|h|^2\bigr)^{\frac p2-1}\ \le\ K(h,e)\ \le\ C_2\bigl(e+|h|^2\bigr)^{\frac p2-1}
\qquad(h\in\R^n,\ e>0).
\]
Granting it, the two structure conditions for $A_e$ are obtained from
\emph{two-point increments} of $A$, with no derivative of $A$ ever integrated.

\emph{\textup{(Q1)} for $A_e$, from the monotonicity of $A$.}  Apply
\noderef{StructureMonotonicityBound} with the parameters $n,p$ and with the
present $r$ in the role of its $r$ --- legitimate since $2\le n$, $1<p$ and
$r>0$ --- and let $C_M>0$ be the constant it produces, depending only on
$n,p,r$.  Apply its conclusion to $\varepsilon=0$, the constants
$\gamma_0>0$ and $\gamma_1\le r\gamma_0$, the set $\R^n$ and the field
$(x,z,h)\mapsto A(h)$, whose \noderef{EllipticStructureQ12} hypothesis is the
hypothesis of the theorem; read at every point as above, it gives
\[
C_M\gamma_0\bigl(|a|^2+|b|^2\bigr)^{\frac p2-1}|a-b|^2
\ \le\ \bigl(A(a)-A(b),\,a-b\bigr)
\qquad(a,b\in\R^n).
\tag{MA}
\]
Note that (MA) carries no derivative and excludes no point: it holds for
\emph{all} $a,b$, the degenerate case $a=b$ reading $0\le0$.

Fix $h,\xi\in\R^n$ with $\xi\ne0$ (for $\xi=0$ both sides of (Q1) of the thesis
vanish) and let $0<t\le1$.  Apply (MA) at
\[
a=h+t\xi+\sqrt e\,z,\qquad b=h-t\xi+\sqrt e\,z ,
\]
so that $a-b=2t\xi$ and, by the parallelogram law
$|u+v|^2+|u-v|^2=2|u|^2+2|v|^2$ with $u=h+\sqrt e\,z$ and $v=t\xi$,
\[
|a|^2+|b|^2=2|h+\sqrt e\,z|^2+2t^2|\xi|^2 .
\]
Thus, for every $z$,
\[
4C_M\gamma_0t^2|\xi|^2
\bigl(2|h+\sqrt e\,z|^2+2t^2|\xi|^2\bigr)^{\frac p2-1}
\ \le\ 2t\bigl(A(a)-A(b),\,\xi\bigr),
\]
and the base $2|h+\sqrt e\,z|^2+2t^2|\xi|^2$ is strictly positive because
$t\ne0$ and $\xi\ne0$, so no reading of $0^{p/2-1}$ is involved.  Both sides are
continuous in $z$ --- the left because its base is positive, the right because
$A$ is continuous on $\R^n$ --- and both are bounded on the compact
$\operatorname{supp}\rho\subseteq\overline{B_1(0)}$, since for $|z|\le1$ the two
points $a,b$ lie in the compact ball
$\overline{B_{|h|+t|\xi|+\sqrt e}(0)}$ on which $A$ is bounded, being
continuous.  So both sides are integrable against $\rho(z)\,dz$ (Preamble,
clause (a) of the domination item, against a constant majorant on a set of
finite measure); integrating and using
\[
\int_{\R^n}\bigl[A(h+t\xi+\sqrt e\,z)-A(h-t\xi+\sqrt e\,z)\bigr]\rho(z)\,dz
=A_e(h+t\xi)-A_e(h-t\xi)
\]
--- the two renormalising constants cancel --- and dividing by $2t>0$ gives
\[
C_M\gamma_0|\xi|^2L(t)
\ \le\ \frac{\bigl(A_e(h+t\xi)-A_e(h-t\xi),\,\xi\bigr)}{2t},
\qquad
L(t)=\int_{\R^n}\bigl(2|h+\sqrt e\,z|^2+2t^2|\xi|^2\bigr)^{\frac p2-1}\rho(z)\,dz .
\tag{Q1-t}
\]
As $t\downarrow0$ the right-hand side of (Q1-t) tends to
$(DA_e(h)\xi,\xi)$: $A_e$ is differentiable at $h$ by the convolution paragraph,
so $\bigl(A_e(h\pm t\xi)-A_e(h)\bigr)/t\to\pm DA_e(h)\xi$, and the symmetric
difference quotient is the half-sum of the two.  For the left-hand side we
bound $L(t)$ below, in each of the two sign cases for $q=\tfrac p2-1$, by a
quantity whose limit is explicit; no interchange of limit and integral is
needed.
\begin{itemize}
\item If $p\ge2$, then $q\ge0$ and
  $2|h+\sqrt e\,z|^2+2t^2|\xi|^2\ge2|h+\sqrt e\,z|^2$, so
  $L(t)\ge2^qK(h,e)\ge2^qc_1(e+|h|^2)^q$ by the sandwich, a bound independent
  of $t$.
\item If $1<p<2$, then $q<0$ and for $|z|\le1$ one has
  $|h+\sqrt e\,z|^2\le2|h|^2+2e=2W$ with $W=e+|h|^2$, so the base is at most
  $4W+2t^2|\xi|^2$ and, the exponent being negative,
  $L(t)\ge\bigl(4W+2t^2|\xi|^2\bigr)^q\int\rho=\bigl(4W+2t^2|\xi|^2\bigr)^q$,
  whose limit as $t\downarrow0$ is $4^qW^q$, since $W\ge e>0$.
\end{itemize}
Writing $c_4=2^qc_1$ when $p\ge2$ and $c_4=4^q$ when $1<p<2$ --- positive in
both cases and depending only on $n$ and $p$ --- and letting $t\downarrow0$ in
(Q1-t),
\[
C_M c_4\,\gamma_0\bigl(e+|h|^2\bigr)^{\frac p2-1}|\xi|^2
\ \le\ \bigl(DA_e(h)\xi,\xi\bigr)
=\sum_{i,j}D_{h_j}A_e^i(h)\,\xi^i\xi^j ,
\]
which is (Q1) of the thesis for $A_e$ at $\varepsilon=e$ with
$\gamma_0^{(e)}=c_0\gamma_0$, where $c_0:=C_Mc_4>0$.  The guard
$e+|h|^2\ne0$ of \noderef{EllipticStructureQ12} is automatic for $e>0$.

\emph{The two Carath\'eodory clauses for $A_e$, and for $A$.}  Both fields are
read as fields on $\R^n\times\R\times\R^n$ that do not depend on the point or on
the value: $(x,z,h)\mapsto A_e(h)$ and $(x,z,h)\mapsto A(h)$.  For fixed $z,h$
the map $x\mapsto A_e(h)$ is therefore constant, hence Lebesgue measurable,
which is the measurability clause of \noderef{EllipticStructureQ12}; and for
fixed $x$ the map $(z,h)\mapsto A_e(h)$ is the composition of $A_e$ with the
continuous projection $(z,h)\mapsto h$, so it is continuous on $\R\times\R^n$ as
soon as $A_e$ is, which the display above records ($A_e\in C^\infty(\R^n)$).
That is the joint-continuity clause.  The same two hold for $A$ itself, with
continuity of $A$ supplied by clause \texttt{hcont} of its own
\noderef{EllipticStructureQ12} hypothesis --- and indeed the hypothesis bundle
for $A$ already carries both Carath\'eodory clauses, so nothing has to be
re-derived for it; they are recorded here only because the \emph{conclusion}
asserts a \noderef{EllipticStructureQ12} bundle for $A_e$ and every clause of
that bundle must be produced.

\emph{\textup{(Q2)} for $A_e$, from segment increments of $A$.}  Fix
$h\in\R^n$ and indices $i,j$, and let $t\ne0$.  First, an increment bound for
$A$ itself.  Fix $b\in\R^n$ and put $\varphi(\varsigma)=A^i(b+\varsigma t e_j)$
for $\varsigma\in[0,1]$.  Then $\varphi$ is continuous on $[0,1]$, because $A$
is continuous on $\R^n$ (clause \texttt{hcont}).  The set
$T=\{\varsigma\in\R:b+\varsigma te_j=0\}$ has at most one element, the map
$\varsigma\mapsto b+\varsigma te_j$ being injective for $t\ne0$.  Off $T$ the
field $A$ is of class $C^1$ (clause \texttt{hdiff} of
\noderef{EllipticStructureQ12} at $\varepsilon=0$, whose domain
$D(0)=\{h':0+|h'|^2\ne0\}$ is $\R^n\setminus\{0\}$), so $\varphi$ is
differentiable there with
$\varphi'(\varsigma)=t\,D_{h_j}A^i(b+\varsigma te_j)$, and (Q2) of the thesis read at
every point gives
\[
|\varphi'(\varsigma)|\ \le\ |t|\,\gamma_1\bigl(|b+\varsigma te_j|^2\bigr)^{\frac p2-1}
\qquad(\varsigma\in[0,1]\setminus T).
\]
That majorant is interval-integrable on $[0,1]$ by
\noderef{SegmentWeightIntegrable}, applied with the present $n,p$, with
$\varepsilon=0$ and with the two endpoints $b+te_j$ and $b$: its conclusion is
the interval integrability on $[0,1]$ of
$\varsigma\mapsto\bigl(0+|\varsigma(b+te_j)+(1-\varsigma)b|^2\bigr)^{p/2-1}$,
and $\varsigma(b+te_j)+(1-\varsigma)b=b+\varsigma te_j$.  Its hypotheses are
$1<p$ and $0\le\varepsilon$, both present.  Hence, by the fundamental theorem of
calculus in interval-integral form (Preamble) applied on $[0,1]$ with the single
exceptional point of $T$ --- if $T\cap[0,1]=\emptyset$ the theorem applies
directly; otherwise write $T\cap[0,1]=\{\varsigma_\ast\}$, apply the theorem on
$[0,\varsigma']$ and $[\varsigma'',1]$ for
$\varsigma'<\varsigma_\ast<\varsigma''$ and let
$\varsigma'\uparrow\varsigma_\ast$, $\varsigma''\downarrow\varsigma_\ast$, both
$\varphi$ and the integral of $\varphi'$ being continuous at $\varsigma_\ast$,
the latter because its integrand is interval-integrable --- we get
\[
\bigl|A^i(b+te_j)-A^i(b)\bigr|
=\bigl|\varphi(1)-\varphi(0)\bigr|
\ \le\ |t|\,\gamma_1\int_0^1\bigl(|b+\varsigma te_j|^2\bigr)^{\frac p2-1}d\varsigma .
\tag{IA}
\]
Now take $b=h+\sqrt e\,z$ in (IA).  Because the two renormalising constants
cancel in the difference,
\[
A_e^i(h+te_j)-A_e^i(h)
=\int_{\R^n}\bigl[A^i(h+te_j+\sqrt e\,z)-A^i(h+\sqrt e\,z)\bigr]\rho(z)\,dz ,
\]
the integrand being continuous in $z$ and bounded on the compact
$\operatorname{supp}\rho$ (as $A$ is continuous on the compact ball
$\overline{B_{|h|+|t|+\sqrt e}(0)}$), hence integrable against $\rho$.  The norm
inequality for the integral (Preamble) and then (IA) pointwise in $z$ give
\[
\bigl|A_e^i(h+te_j)-A_e^i(h)\bigr|
\ \le\ |t|\,\gamma_1\int_{\R^n}\!\!\int_0^1
\bigl(|h+\varsigma te_j+\sqrt e\,z|^2\bigr)^{\frac p2-1}d\varsigma\,\rho(z)\,dz ,
\]
the right-hand side being for the moment a nonnegative integral with values in
$[0,\infty]$.  Its integrand is nonnegative and jointly measurable in
$(z,\varsigma)$: the map $y\mapsto\bigl(|y|^2\bigr)^{p/2-1}$ is continuous on
$\R^n\setminus\{0\}$ and takes the value $0^{p/2-1}$ at the origin, hence is
Borel measurable (Preamble: measurability of superlevel sets and of countable
operations on them), and $(z,\varsigma)\mapsto h+\varsigma te_j+\sqrt e\,z$ is
continuous.  So Tonelli (Preamble) exchanges the two integrations, and the inner
integral becomes $K(h+\varsigma te_j,e)$, which the sandwich bounds by
$C_2\bigl(e+|h+\varsigma te_j|^2\bigr)^{p/2-1}$ --- in particular the
right-hand side is finite.  Therefore
\[
\frac{\bigl|A_e^i(h+te_j)-A_e^i(h)\bigr|}{|t|}
\ \le\ C_2\gamma_1\int_0^1\bigl(e+|h+\varsigma te_j|^2\bigr)^{\frac p2-1}d\varsigma .
\]
Let $t\to0$.  Because $e>0$, the function
$b\mapsto\bigl(e+|b|^2\bigr)^{p/2-1}$ is continuous on $\R^n$, hence uniformly
continuous on the compact ball $\overline{B_1(h)}$, which contains
$h+\varsigma te_j$ for all $\varsigma\in[0,1]$ and $|t|\le1$; so the integrand
converges to $\bigl(e+|h|^2\bigr)^{p/2-1}$ uniformly in $\varsigma$ and the
integral converges to that value.  The left-hand side converges to
$\bigl|D_{h_j}A_e^i(h)\bigr|$, $A_e$ being differentiable at $h$.  Hence
\[
\bigl|D_{h_j}A_e^i(h)\bigr|
\ \le\ C_2\gamma_1\bigl(e+|h|^2\bigr)^{\frac p2-1},
\]
which is (Q2) of the thesis for $A_e$ at $\varepsilon=e$ with
$\gamma_1^{(e)}=C_2\gamma_1$.  Together with the smoothness of $A_e$, which
supplies clauses \texttt{hcont} and \texttt{hdiff} of
\noderef{EllipticStructureQ12} on all of $\R^n$, this completes the bundle for
$A_e$ at $\varepsilon=e$ with the constants
$\gamma_0^{(e)}=c_0\gamma_0$ and $\gamma_1^{(e)}=C_2\gamma_1$.  Their ratio
\[
\frac{\gamma_1^{(e)}}{\gamma_0^{(e)}}
=\frac{C_2}{c_0}\cdot\frac{\gamma_1}{\gamma_0}\ \le\ \frac{C_2}{c_0}\,r=:r'
\]
does not depend on $e$, which is clause (1) with the $r'$ of the statement.

\emph{Proof of the sandwich.}  Write $q=\tfrac p2-1$ and $W=e+|h|^2$.  For
$|z|\le1$ one has $|h+\sqrt e\,z|^2\le2|h|^2+2e\le2W$, which settles the two
cases where the sign of $q$ is convenient: for $p\ge2$ ($q\ge0$) it gives
$K\le2^qW^q$ pointwise in $z$ and hence after integration against $\rho$, since
$\int\rho=1$; and for $1<p<2$ ($q<0$) it gives $K\ge2^qW^q$ the same way.
\emph{One point must be discarded in the second case.}  For $q<0$ the step from
$x\le y$ to $x^q\ge y^q$ needs $x>0$, and $x=|h+\sqrt e\,z|^2$ vanishes at the
single $z=-h/\sqrt e$, where pm-0024's convention makes $x^q=0^q=0$, which is
\emph{not} $\ge(2W)^q>0$.  That $z$ is one point of $\R^n$, hence a Lebesgue
null set, so the pointwise inequality holds almost everywhere and the integrated
inequality $K\ge2^qW^q$ is unaffected (Preamble, clause (a) of the domination
item and the monotonicity of the integral over a fixed set, both insensitive to
a null set).  In the case $p\ge2$ no point need be discarded: there
$0^q=0\le(2W)^q$ for $q>0$, and at $p=2$ both sides are $1$.  The two remaining
cases are the ones that do the work.

\emph{The lower bound for $p\ge2$.}  Here the integrand must be kept away from
$0$, so $|h+\sqrt e\,z|$ has to be bounded below on a set of positive
$\rho$-measure; it cannot be bounded below for every $z$, since $h+\sqrt e\,z$
vanishes at one $z$.  Let $u$ be a unit vector with $(u,h)\ge0$ and put
$Z=B_{1/8}(\tfrac34u)\subseteq B_{7/8}(0)$.  For $z\in Z$ write $z=\tfrac34u+w$
with $|w|\le\tfrac18$; from
\[
\Bigl|\frac h{\sqrt e}+\tfrac34u\Bigr|^2
=\frac{|h|^2}e+\frac32\frac{(u,h)}{\sqrt e}+\frac9{16}\ \ge\ \frac9{16}
\]
we get $\bigl|\tfrac h{\sqrt e}+z\bigr|\ge\tfrac34-\tfrac18=\tfrac58$, i.e.
\[
|h+\sqrt e\,z|\ \ge\ \tfrac58\sqrt e\quad(z\in Z),
\qquad
\int_Z\rho\ \ge\ \kappa\,\bigl|B_{1/8}(0)\bigr|=:m>0 ,
\]
and the mass bound is uniform in $u$ exactly because $\rho\ge\kappa>0$ on
$\overline{B_{7/8}(0)}$, which contains every such $Z$.  Now split.  If
$|h|\le2\sqrt e$ then $W\le5e$, so on $Z$ one has
$|h+\sqrt e\,z|^2\ge\tfrac{25}{64}e\ge\tfrac5{64}W$ and hence
$K\ge m\bigl(\tfrac5{64}\bigr)^qW^q$.  If $|h|>2\sqrt e$ then
$|h+\sqrt e\,z|\ge|h|-\sqrt e\ge\tfrac12|h|$ for \emph{every} $|z|\le1$, while
$W\le\tfrac54|h|^2$, so $|h+\sqrt e\,z|^2\ge\tfrac14|h|^2\ge\tfrac15W$ and
$K\ge\bigl(\tfrac15\bigr)^qW^q$.  Since $m\le1$ and $\tfrac5{64}<\tfrac15$, the
constant $c_1=m\bigl(\tfrac5{64}\bigr)^q$ serves in both cases.

\emph{The upper bound for $1<p<2$.}  Here the integrand is singular where
$h+\sqrt e\,z$ vanishes, and what is needed is that the singularity is
integrable against $\rho$ \emph{uniformly} in $h$ and $e$.  Scaling, with
$a=h/\sqrt e$,
\[
K(h,e)=e^{\frac{p-2}2}I(a),
\qquad
I(a)=\int_{\R^n}\rho(z)\,|a+z|^{p-2}\,dz .
\]
If $|a|\le2$ then $|a+z|\le3$ on $\operatorname{supp}\rho$, so, substituting
$w=a+z$ (Preamble: invariance of Lebesgue measure under translations) and
bounding $\rho$ by $\|\rho\|_\infty$,
\[
I(a)\ \le\ \|\rho\|_\infty\int_{\{|w|\le3\}}|w|^{p-2}\,dw ,
\]
and that integral is finite with a bound depending only on $n$ and $p$.
\emph{We do not use the polar-coordinate closed form} $n\omega_n\,3^{n+p-2}/(n+p-2)$,
which no node states and which the Preamble does not declare; dyadic annuli and
the declared volume-of-a-ball item suffice, and only finiteness is needed.

Write $s=p-2$, so $s<0$ and $n+s=n+p-2>0$, since $n\ge2$ and $p>1$.  The value
of the integrand at $w=0$ is $0^{s}=0$ by pm-0024's convention, so the origin
contributes nothing and we may integrate over $\{0<|w|\le3\}$ instead.  For each
integer $k\ge0$ put
\[
A_k=\bigl\{w:3\cdot2^{-(k+1)}<|w|\le3\cdot2^{-k}\bigr\} ,
\]
a disjoint family whose union is exactly $\{0<|w|\le3\}$: given $w$ with
$0<|w|\le3$ there is exactly one $k\ge0$ with
$3\cdot2^{-(k+1)}<|w|\le3\cdot2^{-k}$.  On $A_k$ the exponent being negative
gives $|w|^{s}\le\bigl(3\cdot2^{-(k+1)}\bigr)^{s}$, and
$A_k\subseteq B_{3\cdot2^{-k}+\eta}(0)$ for every $\eta>0$, so
$|A_k|\le\omega_n\bigl(3\cdot2^{-k}+\eta\bigr)^n$ (Preamble: the volume of a
ball); letting $\eta\downarrow0$ gives
$|A_k|\le\omega_n\bigl(3\cdot2^{-k}\bigr)^n$, with no need to measure the
bounding sphere.
Hence
\[
\int_{A_k}|w|^{s}\,dw
\ \le\ \omega_n3^{n}2^{-kn}\cdot3^{s}2^{-(k+1)s}
=\omega_n3^{\,n+s}2^{-s}\bigl(2^{-(n+s)}\bigr)^{k} .
\]
Since $n+s>0$ we have $2^{-(n+s)}<1$, so these terms are a convergent geometric
series with sum
\[
S:=\omega_n3^{\,n+s}2^{-s}\bigl(1-2^{-(n+s)}\bigr)^{-1}<\infty .
\]
To conclude, put $f=|w|^{s}\mathbf1_{\{0<|w|\le3\}}$ and, for each integer $K\ge0$,
$f_K=|w|^{s}\mathbf1_{E_K}$ with $E_K=\bigcup_{k\le K}A_k
=\{3\cdot2^{-(K+1)}<|w|\le3\}$.  Each $f_K$ is nonnegative and measurable,
$f_K\le f_{K+1}$, and $f_K\to f$ pointwise, since every $w$ with $0<|w|\le3$
lies in $E_K$ for all large $K$.  By finite additivity of the integral over the
disjoint $A_0,\dots,A_K$, $\int f_K\le\sum_{k=0}^{K}\omega_n3^{n+s}2^{-s}
\bigl(2^{-(n+s)}\bigr)^{k}\le S$ for every $K$; and Fatou's lemma (Preamble),
applied to the sequence $(f_K)$ whose pointwise limit --- and hence pointwise
$\liminf$ --- is $f$, gives
\[
\int_{\{0<|w|\le3\}}|w|^{s}\,dw=\int f\ \le\ \liminf_K\int f_K\ \le\ S .
\]
So $I(a)\le\|\rho\|_\infty S<\infty$ when $|a|\le2$.  This, with
$\|\rho\|_\infty$ and $\omega_n$, is where the dimension enters, and the
condition $n+p-2>0$ is exactly what makes the geometric ratio less than $1$.  If $|a|\ge2$ then $|a+z|\ge|a|-1\ge\tfrac12|a|$ on
$\operatorname{supp}\rho$, so $I(a)\le2^{2-p}|a|^{p-2}$.  As $1+|a|\le\tfrac32|a|$
in the second case and $1+|a|\le3$ in the first, both give
$I(a)\le C(n,p)\,(1+|a|)^{p-2}$, whence
\[
K(h,e)\ \le\ C\,e^{\frac{p-2}2}\Bigl(1+\frac{|h|}{\sqrt e}\Bigr)^{p-2}
=C\,\bigl(\sqrt e+|h|\bigr)^{p-2}\ \le\ C\,W^{\frac{p-2}2}=C\,W^q ,
\]
the last step because $(\sqrt e+|h|)^2\ge e+|h|^2=W$ and $p-2<0$.  This proves
the sandwich and with it clause (1), for every $p>1$ at once.

\emph{Two remarks on the family.}  First, the sandwich constants depend on the
dimension --- through $\|\rho\|_\infty$, $\kappa$, $|B_{1/8}(0)|$ and the
volume constant $\omega_n=|B_1(0)|$ of the Preamble --- so the number it
produces is
$r'=r'(n,p,\gamma_1/\gamma_0)$ and not a function of $p$ and the ratio alone;
this is why the statement records the dependence as $r'(n,p,r)$ and binds $r'$
after $n$.  Nothing is lost by this: both consumers bind $n$ first and their own
constants already depend on $n$.  Second, the additive family
$A_e(h)=A(h)+c\gamma_0e^{p/2-1}h$, to which the counterexample in the statement
block refers, also serves for $p\ge2$, but not for $1<p<2$, where the exponent
$\tfrac p2-1$ is negative; mollification is a single family for every $p>1$ and
is what is used here.

\emph{What the hypotheses contain when $1<p<2$.}  One fact about the class being
approximated should be recorded, because the $1<p<2$ case is the one this node
exists for.  Read with differentiability demanded on all of $\R^n$, the
structure conditions at $\varepsilon=0$ would be met by nothing at all when
$1<p<2$: (Q1) of the thesis forces $|DA(h)|\ge\gamma_0|h|^{p-2}\to\infty$ as
$h\to0$, while $C^1$ on all of $\R^n$ makes $DA$ continuous at $0$, and no field
does both --- an argument that uses only $h\ne0$ and is therefore independent of
how $0^{p/2-1}$ is read.  That is why \noderef{EllipticStructureQ12} asks for
continuity on $\R^n$ together with $C^1$ regularity and
(Q1) of the thesis--(Q2) of the thesis on $D(\varepsilon)=\{h:\varepsilon+|h|^2\ne0\}$
only, which at $\varepsilon=0$ is $\R^n\setminus\{0\}$; with that reading the
model field $A(h)=|h|^{p-2}h$ --- the one this whole section is about --- lies in
the class, with $\gamma_0=p-1$ and $\gamma_1=3-p$.  The proof above is written to
use exactly what that reading supplies: (Q1) of the thesis--(Q2) of the thesis at points
$h\ne0$, the continuity of $A$ on all of $\R^n$, and the local integrability of
$|DA|$.  For the two regularized problems nothing is lost either way, since
$D(e)=\R^n$ for every $e>0$.

\medskip
We now prove clauses (2)--(5).  The argument is in nine steps: the structure
conditions are first read pointwise and the regularized solutions produced from
(C2) (Steps 1--3), then the two energy identities are established by
approximating the test function $u_e-u$ (Step 4), then the gradients are bounded
uniformly (Step 5) and their difference estimated (Steps 6--8), and finally the
sequence of the conclusion is chosen (Step 9).

\emph{Step 1: the structure constants are positive.}  The promotion of
(Q1) of the thesis and (Q2) of the thesis from ``for a.e.\ $x$'' to ``for every $x$'' was
carried out at the head of this proof, before the family was constructed,
because the construction already needs it; it applies verbatim to $A_e$, whose
\noderef{EllipticStructureQ12} bundle at $\varepsilon=e$ was established there
and which likewise does not depend on $x$.  We use both clauses as statements
about the fields of $h$ alone from here on.

In particular $\gamma_1>0$: take $h=e_1$, so that $0+|h|^2=1\ne0$, and
$\xi=e_1$ in (Q1) of the thesis and $i=j=1$ in (Q2) of the thesis, obtaining
\[
\gamma_0=\gamma_0\cdot1^{\frac p2-1}|e_1|^2
\ \le\ \sum_{i,j}D_{h_j}A^i(e_1)e_1^ie_1^j=D_{h_1}A^1(e_1)
\ \le\ \bigl|D_{h_1}A^1(e_1)\bigr|\ \le\ \gamma_1\cdot1^{\frac p2-1}=\gamma_1 ,
\]
so $0<\gamma_0\le\gamma_1$.  The same argument applied to $A_e$ gives
$0<c_0\gamma_0\le C_2\gamma_1$; both structure constants of $A_e$ are therefore
positive, which is what the cited results require.

\emph{Step 2: two quantitative consequences of the structure conditions.}  Apply
\noderef{StructureMonotonicityBound} with the parameters $n,p$ and with $r'$ in
the role of its $r$ --- legitimate since $2\le n$, $1<p$ and $r'>0$ --- and let
$C_M'>0$ be the constant it produces.  (This is a second application of that
node; the first, with the present $r$ and the field $A$, produced the $C_M$ of
the family paragraph.  The two constants are different numbers and both depend
only on $n,p$ and the respective ratio bound.)  Applied to the field
$(x,z,h)\mapsto A_e(h)$ on $\R^n$ with $\varepsilon=e$ and the constants
$c_0\gamma_0>0$, $C_2\gamma_1\le r'c_0\gamma_0$ of clause (1), and read at every
point as in Step 1, it gives
\[
C_M'c_0\gamma_0\bigl(e+|a|^2+|b|^2\bigr)^{\frac p2-1}|a-b|^2
\ \le\ \bigl(A_e(a)-A_e(b),\,a-b\bigr)
\qquad(a,b\in\R^n,\ e>0).
\tag{M}
\]
Write $c_3=C_M'c_0\gamma_0>0$.

Apply \noderef{StructureGrowthBound} with the parameters $n,p$ --- its
hypotheses are $2\le n$ and $1<p$ --- and let $c_G>0$ be the constant it
produces.  Its conclusion, like that of \noderef{StructureMonotonicityBound}, is
stated ``for a.e.\ $x$'' with an $x$-free right-hand side, so it is read at
every point exactly as at the head of this proof.  Applied to $A_e$ (with
$\varepsilon=e$, $c_0\gamma_0>0$, $C_2\gamma_1>0$) and to $A$ (with
$\varepsilon=0$, $\gamma_0>0$, $\gamma_1>0$), and using $A_e(0)=0$ and $A(0)=0$,
it gives
\[
|A_e(h)|\le c_GC_2\gamma_1\bigl(e+|h|^2\bigr)^{\frac p2-1}|h| ,
\qquad
|A(h)|\le c_G\gamma_1\bigl(|h|^2\bigr)^{\frac p2-1}|h|=c_G\gamma_1|h|^{p-1} .
\]
We record the consequence used repeatedly: for $0<e\le1$ and every $h$,
\[
|A_e(h)|\ \le\ \Gamma\bigl(1+|h|^2\bigr)^{\frac{p-1}2},
\qquad
|A(h)|\ \le\ \Gamma\bigl(1+|h|^2\bigr)^{\frac{p-1}2},
\qquad
\Gamma:=c_G\gamma_1\max\{1,C_2\} .
\tag{G}
\]
Indeed, if $p\ge2$ then $(e+|h|^2)^{\frac p2-1}\le(1+|h|^2)^{\frac p2-1}$ for
$e\le1$ and $|h|\le(1+|h|^2)^{1/2}$, so the product is at most
$(1+|h|^2)^{\frac{p-1}2}$; if $1<p<2$ the exponent is negative, so
$(e+|h|^2)^{\frac p2-1}\le(|h|^2)^{\frac p2-1}=|h|^{p-2}$ for $h\ne0$, giving
$|A_e(h)|\le c_GC_2\gamma_1|h|^{p-1}\le c_GC_2\gamma_1(1+|h|^2)^{\frac{p-1}2}$,
while at $h=0$ the left-hand side is $0$.  The bound for $A$ is the displayed one
together with $|h|^{p-1}\le(1+|h|^2)^{\frac{p-1}2}$.

\emph{Step 3: clauses (2) and (3).}  Fix $e$ with $0<e\le1$.  The hypothesis of
the present theorem assumes the cited solvability result (C2)
(\noderef{CitedDirichletW1p}) for \emph{every} field and all parameters with
positive structure constants; instantiate it at $\varepsilon=e$, at the
constants $c_0\gamma_0>0$ and $C_2\gamma_1>0$ of Step 1, and at the field $A_e$.
The hypotheses of (C2) itself hold: $A_e$ satisfies
(Q1) of the thesis--(Q2) of the thesis on $\R^n$ with those parameters by clause (1),
$0\le e$, and $A_e(0)=0$ by construction.  Applied on the ball $B_R(x_0)$, with
$R>0$, and with the boundary datum the pair $(u,g)$, which satisfies
$u\in W^{1,p}(B_R(x_0))$ with weak gradient $g$ (\noderef{MemW1pOn}) by
hypothesis, it produces $u_e$ and $g_e$ with
\[
u_e\ \text{a weak solution of }\operatorname{div}A_e(\nabla u_e)=0\ \text{in }B_R(x_0)
\ (\noderef{IsWeakSolutionOn}),
\tag{2$_e$}
\]
\[
u_e-u\in W^{1,p}_0\bigl(B_R(x_0)\bigr)\ \text{with weak gradient }g_e-g
\ (\noderef{MemW1pZeroOn}).
\tag{3$_e$}
\]
These are clauses (2) and (3) at the index in question.  \emph{The
integrability clause of \noderef{IsWeakSolutionOn} comes with
$(2_e)$ and is not re-derived here}: that clause is part of the predicate the
cited result (C2) (\noderef{CitedDirichletW1p}) asserts of the solution it
produces, so it is handed over together with the integral identity, for every
admissible test pair, at no cost to this proof.  Note that
$(3_e)$ contains, as its first clause, $u_e-u\in W^{1,p}(B_R(x_0))$ with weak
gradient $g_e-g$; in particular $|g_e-g|\in L^p(B_R(x_0))$ and hence, with
$|g|\in L^p(B_R(x_0))$ and the $L^p$ triangle inequality (Minkowski, Preamble),
$|g_e|\in L^p(B_R(x_0))$.  Write
\[
Y_e=\int_{B_R(x_0)}|g_e|^p\,dx<\infty,
\qquad
K=\int_{B_R(x_0)}|g|^p\,dx<\infty,
\qquad
X_e=\Bigl(\int_{B_R(x_0)}|g_e-g|^p\,dx\Bigr)^{1/p}.
\]

\emph{Step 4: the two energy identities.}  We claim
\[
\int_{B_R(x_0)}\bigl(A_e(g_e),\,g_e-g\bigr)dx=0,
\qquad
\int_{B_R(x_0)}\bigl(A(g),\,g_e-g\bigr)dx=0 .
\tag{E}
\]
The test class of \noderef{IsWeakSolutionOn} consists of pairs $(\varphi,g_\varphi)$
with $\varphi\in W^{1,p}(B_R(x_0))$ of weak gradient $g_\varphi$, $\varphi$ of
compact support and $\operatorname{supp}\varphi\subseteq B_R(x_0)$; the function
$u_e-u$ is not of compact support, so we approximate.  By
\noderef{MemW1pZeroOn} the membership $(3_e)$ provides functions
$\varphi_k\in C^\infty(\R^n)$, $k\in\mathbb{N}$, each with compact support and
$\operatorname{supp}\varphi_k\subseteq B_R(x_0)$, such that
\[
\int_{B_R(x_0)}\bigl|(u_e-u)-\varphi_k\bigr|^p\,dx\to0,
\qquad
\int_{B_R(x_0)}\bigl|(g_e-g)-\nabla\varphi_k\bigr|^p\,dx\to0
\tag{4a}
\]
as $k\to\infty$ (the second integral is the one written in that node with the
sum over the components under the power $p/2$, which is
$|(g_e-g)-\nabla\varphi_k|^p$).  Each pair $(\varphi_k,\nabla\varphi_k)$ is
admissible: $\varphi_k\in W^{1,p}(B_R(x_0))$ with weak gradient
$\nabla\varphi_k$ by \noderef{SmoothCompactSupportMemW1p}, whose hypotheses
$1\le n$, $1\le p$, openness of $B_R(x_0)$, smoothness of $\varphi_k$ and
compactness of its support inside the ball all hold.  So $(2_e)$ and the
hypothesis that $u$ is a weak solution of $\operatorname{div}A(\nabla u)=0$ in
$B_R(x_0)$ give, for every $k$,
\[
\int_{B_R(x_0)}\bigl(A_e(g_e),\nabla\varphi_k\bigr)dx=0,
\qquad
\int_{B_R(x_0)}\bigl(A(g),\nabla\varphi_k\bigr)dx=0 ,
\]
the vanishing right-hand side of the equations making the $B$-term absent.
Now let $k\to\infty$.  By (G) and $(p-1)p'=p$, where $p'=p/(p-1)$,
\[
\int_{B_R(x_0)}|A_e(g_e)|^{p'}dx
\le\Gamma^{p'}\int_{B_R(x_0)}\bigl(1+|g_e|^2\bigr)^{\frac p2}dx
\le\Gamma^{p'}2^{\frac p2}\bigl(|B_R|+Y_e\bigr)<\infty ,
\]
using $(1+t^2)^{p/2}\le2^{p/2}(1+t^p)$ for $t\ge0$; so $|A_e(g_e)|\in L^{p'}$,
and likewise $|A(g)|\in L^{p'}$ with $Y_e$ replaced by $K$.  Hölder's
inequality (Preamble) therefore gives
\[
\Bigl|\int_{B_R(x_0)}\bigl(A_e(g_e),\nabla\varphi_k-(g_e-g)\bigr)dx\Bigr|
\le\bigl\|\,|A_e(g_e)|\,\bigr\|_{L^{p'}}
\bigl\|\,|\nabla\varphi_k-(g_e-g)|\,\bigr\|_{L^{p}}\longrightarrow0
\]
by $(4a)$, and the same with $A(g)$ in place of $A_e(g_e)$.  Hence both
integrals in (E) are limits of integrals that vanish for every $k$, which proves
(E).

\emph{Step 5: a bound on $Y_e$ uniform in $e$.}  The first identity in (E) reads
$\int(A_e(g_e),g_e)=\int(A_e(g_e),g)$.  For the left-hand side use (M) with
$a=g_e(x)$, $b=0$ and $A_e(0)=0$:
\[
\bigl(A_e(g_e(x)),g_e(x)\bigr)\ \ge\ c_3\bigl(e+|g_e(x)|^2\bigr)^{\frac p2-1}|g_e(x)|^2 .
\]
If $p\ge2$ the exponent is nonnegative, so
$(e+|h|^2)^{\frac p2-1}\ge(|h|^2)^{\frac p2-1}=|h|^{p-2}$ and the right-hand
side is at least $c_3|h|^p$ with $h=g_e(x)$.  If $1<p<2$ the exponent is
negative; then at points where $|h|^2\ge e$ we have $e+|h|^2\le2|h|^2$ and hence
$(e+|h|^2)^{\frac p2-1}\ge2^{\frac p2-1}|h|^{p-2}$, so the right-hand side is at
least $c_32^{\frac p2-1}|h|^p$ there, while at the remaining points it is
nonnegative and $|h|^p<e^{p/2}$.  With $\kappa_p=\min\{1,2^{\frac p2-1}\}>0$,
integrating and using monotonicity of the integral (Preamble),
\[
\int_{B_R(x_0)}\bigl(A_e(g_e),g_e\bigr)dx
\ \ge\ c_3\kappa_p\Bigl(Y_e-e^{\frac p2}|B_R|\Bigr).
\]
For the right-hand side, Hölder's inequality and the computation of Step 4 give
\[
\Bigl|\int_{B_R(x_0)}\bigl(A_e(g_e),g\bigr)dx\Bigr|
\le\Gamma2^{\frac p{2p'}}\bigl(|B_R|+Y_e\bigr)^{1/p'}K^{1/p}
=:b_0\bigl(|B_R|+Y_e\bigr)^{\frac{p-1}p},
\]
with $b_0=\Gamma2^{\frac p{2p'}}K^{1/p}$ independent of $e$.  Putting
$S=|B_R|+Y_e$ and using $e\le1$,
\[
c_3\kappa_pS\ \le\ c_3\kappa_p\bigl(|B_R|+e^{\frac p2}|B_R|\bigr)+b_0S^{\theta}
\ \le\ 2c_3\kappa_p|B_R|+b_0S^{\theta},
\qquad \theta=\frac{p-1}p\in(0,1),
\]
that is, $S\le a_0+b_1S^{\theta}$ with $a_0=2|B_R|$ and
$b_1=b_0/(c_3\kappa_p)$.  If $S>2a_0$ then $S/2\le b_1S^{\theta}$, so
$S^{1-\theta}\le2b_1$ and $S\le(2b_1)^{1/(1-\theta)}$.  Hence
\[
|B_R|+Y_e\ \le\ E:=\max\bigl\{2a_0,\ (2b_1)^{1/(1-\theta)}\bigr\}
\qquad(0<e\le1),
\tag{5a}
\]
a bound independent of $e$.  Consequently, using
$(s+t+w)^{p/2}\le3^{p/2}(s^{p/2}+t^{p/2}+w^{p/2})$ for $s,t,w\ge0$,
\[
\int_{B_R(x_0)}\bigl(e+|g_e|^2+|g|^2\bigr)^{\frac p2}dx
\ \le\ 3^{\frac p2}\bigl(|B_R|+Y_e+K\bigr)\ \le\ E_1:=3^{\frac p2}(E+K)
\qquad(0<e\le1).
\tag{5b}
\]

\emph{Step 6: the difference estimate.}  Subtracting the two identities of (E),
\[
\int_{B_R(x_0)}\bigl(A_e(g_e)-A_e(g),\,g_e-g\bigr)dx
=\int_{B_R(x_0)}\bigl(A(g)-A_e(g),\,g_e-g\bigr)dx ,
\]
because the first identity makes $\int(A_e(g_e),g_e-g)=0$ and the second makes
$\int(A(g),g_e-g)=0$, so both sides equal $-\int(A_e(g),g_e-g)$; all four
integrals are finite by Step 4 and Hölder.  By (M) with $a=g_e(x)$, $b=g(x)$ the
left-hand side is at least $c_3I_e$, where
\[
I_e=\int_{B_R(x_0)}\bigl(e+|g_e|^2+|g|^2\bigr)^{\frac p2-1}|g_e-g|^2\,dx\ \ge\ 0 ,
\]
and by Hölder's inequality the right-hand side is at most $\delta_eX_e$, where
\[
\delta_e=\Bigl(\int_{B_R(x_0)}\bigl|A(g)-A_e(g)\bigr|^{p'}dx\Bigr)^{1/p'} .
\]
Hence
\[
c_3I_e\ \le\ \delta_eX_e\qquad(0<e\le1).
\tag{6}
\]

\emph{Step 7: $\delta_e\to0$ along any sequence $e\downarrow0$.}  First,
$A_e(h)\to A(h)$ for every fixed $h$ as $e\downarrow0$: $A$ is continuous on
$\R^n$ (the clause \texttt{hcont} of \noderef{EllipticStructureQ12}, which is
pointwise in $x$ and therefore a statement about $A$), so
$\sup_{|y-h|\le\sqrt e}|A(y)-A(h)|\to0$ as $e\downarrow0$ by continuity at $h$;
since $\operatorname{supp}\rho\subseteq\overline{B_1(0)}$ and $\int\rho=1$, this
gives $\int A(h+\sqrt e\,z)\rho(z)dz\to A(h)$ and, at $h=0$,
$\int A(\sqrt e\,z)\rho(z)dz\to A(0)=0$; subtracting the two and recalling the
definition of $A_e$ gives the claim.  Second, by (G),
\[
\bigl|A(g(x))-A_e(g(x))\bigr|^{p'}\ \le\ (2\Gamma)^{p'}\bigl(1+|g(x)|^2\bigr)^{\frac p2}
\qquad(0<e\le1),
\]
and the majorant is integrable on $B_R(x_0)$, being at most
$2^{p/2}(2\Gamma)^{p'}(1+|g|^p)$ with $|B_R|<\infty$ and $|g|\in L^p$.  So for
any sequence $e_m\downarrow0$ in $(0,1]$, dominated convergence (Preamble)
applied to the sequence of functions $x\mapsto|A(g(x))-A_{e_m}(g(x))|^{p'}$ ---
which converges pointwise to $0$ by the first part --- gives
$\delta_{e_m}\to0$.

\emph{Step 8: $X_e\le\Psi(\delta_e)$ with $\Psi(\delta)\to0$ as $\delta\to0$.}
If $X_e=0$ there is nothing to prove, so assume $X_e>0$.

\emph{Case $p\ge2$.}  Pointwise $|a-b|^2\le2|a|^2+2|b|^2\le2(e+|a|^2+|b|^2)$,
so $e+|a|^2+|b|^2\ge\frac12|a-b|^2$ and, the exponent $\frac p2-1$ being
nonnegative,
$\bigl(e+|a|^2+|b|^2\bigr)^{\frac p2-1}\ge2^{1-\frac p2}|a-b|^{p-2}$.  Hence
$I_e\ge2^{1-\frac p2}X_e^{\,p}$, and $(6)$ gives
$c_32^{1-\frac p2}X_e^{\,p}\le\delta_eX_e$, i.e.
\[
X_e\ \le\ \Bigl(\frac{2^{\frac p2-1}}{c_3}\,\delta_e\Bigr)^{\frac1{p-1}} .
\]

\emph{Case $1<p<2$.}  Since $e>0$, the quantity
$\Theta(x)=\bigl(e+|g_e(x)|^2+|g(x)|^2\bigr)^{\frac p2-1}$ is finite and
strictly positive, so we may write
$|g_e-g|^p=\bigl(\Theta|g_e-g|^2\bigr)^{\frac p2}\Theta^{-\frac p2}$ pointwise.
Hölder's inequality with the conjugate exponents $\frac2p$ and $\frac2{2-p}$
(Preamble) gives
\[
X_e^{\,p}=\int_{B_R(x_0)}\bigl(\Theta|g_e-g|^2\bigr)^{\frac p2}\Theta^{-\frac p2}dx
\ \le\ I_e^{\frac p2}
\Bigl(\int_{B_R(x_0)}\Theta^{-\frac p{2-p}}dx\Bigr)^{\frac{2-p}2} .
\]
Now $-\frac p{2-p}\bigl(\frac p2-1\bigr)=\frac p{2-p}\cdot\frac{2-p}2=\frac p2$,
so $\Theta^{-\frac p{2-p}}=\bigl(e+|g_e|^2+|g|^2\bigr)^{\frac p2}$ and the last
integral is at most $E_1$ by $(5b)$.  With $(6)$,
\[
X_e^{\,p}\ \le\ \Bigl(\frac{\delta_eX_e}{c_3}\Bigr)^{\frac p2}E_1^{\frac{2-p}2},
\qquad\text{hence}\qquad
X_e^{\frac p2}\le\Bigl(\frac{\delta_e}{c_3}\Bigr)^{\frac p2}E_1^{\frac{2-p}2},
\qquad\text{i.e.}\qquad
X_e\ \le\ \frac{\delta_e}{c_3}\,E_1^{\frac{2-p}p} .
\]
In both cases $X_e$ is bounded by a continuous increasing function of
$\delta_e$ that vanishes at $\delta_e=0$, and the bound is uniform in
$e\in(0,1]$.

\emph{Step 9: clauses (4) and (5), and the choice of the sequence.}  Put
$\tilde e_m=\frac1{m+1}\in(0,1]$ for $m\in\mathbb{N}$, so
$\tilde e_m\downarrow0$, and let $A_{\tilde e_m}$ be the corresponding members of
the family and $u_{\tilde e_m},g_{\tilde e_m}$ the solutions furnished by Step 3.
By Steps 7 and 8,
\[
\int_{B_R(x_0)}\bigl|g_{\tilde e_m}-g\bigr|^p\,dx=X_{\tilde e_m}^{\,p}
\longrightarrow0\qquad(m\to\infty).
\tag{9a}
\]
By the Preamble item on extracting an almost everywhere convergent subsequence
out of $L^q$ convergence, applied at the exponent $p$ on the measure space
$B_R(x_0)$ to the functions $g_{\tilde e_m}-g$, there is a strictly increasing
sequence of indices $m_0<m_1<\cdots$ such that
\[
g_{\tilde e_{m_j}}(x)\longrightarrow g(x)\qquad\text{for almost every }x\in B_R(x_0).
\tag{9b}
\]
Define the data of the conclusion to be that subsequence:
\[
e_j=\tilde e_{m_j},\qquad A^{(j)}=A_{\tilde e_{m_j}},\qquad
u^{(j)}=u_{\tilde e_{m_j}},\qquad g^{(j)}=g_{\tilde e_{m_j}}
\qquad(j\in\mathbb{N}).
\]
All five clauses now hold: $e_j>0$ and $e_j\to0$, because a subsequence of a
sequence tending to $0$ tends to $0$; clause (1) holds at every index, by the
first part of this proof applied at $e=e_j$; clauses (2) and (3) hold at every
index by Step 3; clause (4) holds because a subsequence of the null sequence
$(9a)$ is null; and clause (5) is $(9b)$.

This is the point at which the sequential form of the statement is essential.
What $L^p$ convergence supplies is almost everywhere convergence along
\emph{some} subsequence, and nothing more: for a family indexed by a continuum
of parameters $e$ there is no almost everywhere convergence as $e\downarrow0$ to
be had from $(9a)$.  The conclusion is therefore stated, and is proved here, for
one sequence that the proof itself chooses --- the subsequence extracted in
$(9b)$ --- and the consumers use it in that form.

\emph{How the limit is used.}  Because the ratios
$\gamma_1^{(m)}/\gamma_0^{(m)}$ are bounded by the single number $r'$, any
estimate for the nondegenerate problems whose constant depends only on $n$, $p$
and that ratio holds for every $m$ with one and the same constant, and therefore
passes to the limit and holds for $u$ with $\varepsilon=0$.  This is how
Theorems~Theorem~1 of the thesis and~Theorem~2 of the thesis are reduced to the case $\varepsilon>0$:
each consumer instantiates its own ratio hypothesis $\gamma_1\le r\gamma_0$ at
$r'$.
\end{proof}
